\documentclass[10pt,a4paper]{article}

\usepackage[margin=1.55cm,headheight=14pt,footskip=21pt]{geometry}
\usepackage{amsmath,amssymb,mathtools}
\usepackage{booktabs,array,tabularx}
\usepackage{enumitem}
\usepackage{xcolor}
\usepackage{listings}
\usepackage{microtype}
\usepackage{fancyhdr}
\usepackage[hidelinks]{hyperref}

\definecolor{navy}{RGB}{21,48,82}
\definecolor{lightblue}{RGB}{235,242,250}
\definecolor{codegreen}{RGB}{20,110,70}
\definecolor{codegray}{RGB}{90,90,90}

\hypersetup{
  pdftitle={An Adaptive Market-Making Model for Binary Event Contracts},
  pdfauthor={Chuanjie Wu},
  pdfsubject={Quantitative market making, RFQ and FOK execution, model risk and solvency}
}

\pagestyle{fancy}
\fancyhf{}
\lhead{Adaptive Binary-Option Market Maker}
\rhead{Technical Research Report}
\cfoot{\thepage}
\renewcommand{\headrulewidth}{0.4pt}

\setlength{\parindent}{0pt}
\setlength{\parskip}{2.5pt}
\setlist[itemize]{leftmargin=1.5em,itemsep=1pt,topsep=1.5pt}
\setlist[enumerate]{leftmargin=1.7em,itemsep=1pt,topsep=1.5pt}
\renewcommand{\arraystretch}{1.08}
\raggedbottom

\lstset{
  language=Python,
  basicstyle=\ttfamily\scriptsize,
  keywordstyle=\color{navy}\bfseries,
  commentstyle=\color{codegreen},
  stringstyle=\color{codegray},
  backgroundcolor=\color{lightblue},
  frame=single,
  rulecolor=\color{navy!35},
  breaklines=true,
  showstringspaces=false,
  columns=fullflexible,
  aboveskip=4pt,
  belowskip=4pt
}

\newcommand{\clip}{\operatorname{clip}}
\newcommand{\E}{\mathbb{E}}
\newcommand{\Var}{\operatorname{Var}}
\newcommand{\Cov}{\operatorname{Cov}}
\newcommand{\1}{\mathbf{1}}
\newcommand{\sectionrule}{\vspace{-2pt}\noindent\color{navy}\rule{\textwidth}{0.7pt}\color{black}\vspace{2pt}}

\begin{document}
\hypersetup{pageanchor=false}

% -----------------------------------------------------------------------------
% Page 1: title
% -----------------------------------------------------------------------------
\begin{titlepage}
\thispagestyle{empty}
\centering
\vspace*{1.1cm}
{\color{navy}\rule{\textwidth}{1.2pt}}\\[1.0cm]
{\LARGE\bfseries An Adaptive Market-Making Model\\[4pt]
for Binary Event Contracts\par}
\vspace{0.55cm}
{\large RFQ Quoting, FOK Selection, Portfolio Risk,\\
Adverse Selection and Runtime Engineering\par}
\vspace{0.8cm}
{\color{navy}\rule{0.78\textwidth}{0.8pt}}\\[0.9cm]
{\large Chuanjie Wu\par}
\vspace{0.25cm}
{\normalsize Quantitative Research Technical Report\par}
\vspace{0.2cm}
{\normalsize 29 August 2026\par}

\vfill
\begin{minipage}{0.88\textwidth}
\small
\textbf{Abstract.}
This report develops a research-grade market-making framework for binary event
contracts whose terminal payoff is either zero or one.  The market maker trades
contracts written on a discrete policy rate (FED) and two correlated private-company
valuations (AJR and THR).  The framework combines a fitted stochastic factor model,
analytic and Monte Carlo probability estimation, inventory-aware reservation prices,
adaptive RFQ spreads, adverse-selection diagnostics, conservative FOK acceptance,
and exact maximum-loss cash accounting.  The design has two ordered objectives:
avoid bankruptcy under every admissible execution sequence, then maximise realised
PnL and competitive rank.  Particular attention is given to the separation between
RFQs, where the dealer controls price and displayed size but does not observe the
customer's direction, and FOKs, where direction, price and total size are revealed but
the dealer can only accept or reject.  A state-level caching architecture reduces
runtime without changing the economic logic.  The report states assumptions,
derives the decision rules, maps each rule to implementable Python components, and
specifies an out-of-sample validation and model-governance protocol.

\vspace{0.5cm}
\textbf{Keywords:} binary option; event contract; market making; RFQ; fill-or-kill;
inventory risk; adverse selection; Monte Carlo; solvency; model governance.
\end{minipage}

\vfill
{\footnotesize Research prototype for controlled simulation.  It is not investment advice
and is not represented as a production trading system.}
\end{titlepage}
\hypersetup{pageanchor=true}
\setcounter{page}{2}
\small
\setcounter{tocdepth}{1}

% -----------------------------------------------------------------------------
% Page 2: contents and executive summary
% -----------------------------------------------------------------------------
\thispagestyle{fancy}
\tableofcontents
\vspace{0.2cm}
\section*{Executive summary}
\addcontentsline{toc}{section}{Executive summary}
\sectionrule

The system prices an option by its estimated probability of paying one.  It then
separates the trading decision into a \emph{directional centre} and a \emph{risk
width}.  The directional centre is the reservation price
\[
 r=p-\gamma_t C+\theta F,
\]
where $p$ is fair value, $C$ is covariance with the existing portfolio,
$\gamma_t$ is dynamic risk aversion, and $\theta F$ is a learned order-flow
adjustment.  For an RFQ, the dealer publishes
\[
 b=\lfloor r-h\rfloor_{0.01},\qquad
 a=\lceil r+h\rceil_{0.01},
\]
where $h$ aggregates base compensation, estimation uncertainty, binary payoff
variance, tenor, live volatility, counterparty toxicity and repeat-request risk.
Displayed quantities are bounded by exact cash capacity and concentration limits.

For a FOK, the customer supplies side, price and total quantity.  The dealer accepts
only when apparent edge exceeds model, size, toxicity, concentration and short-tenor
penalties, and when the complete order is solvent under its worst payoff.  Large,
short-dated FOKs are treated as potentially informed, while ordinary central-probability
RFQs may be quoted more competitively.  This separation is the central strategic
contribution: \emph{compete for ordinary flow without applying the same aggressiveness
to orders that reveal unusually strong information}.

\begin{table}[h]
\centering\small
\begin{tabularx}{\textwidth}{>{\bfseries}p{2.7cm}X X}
\toprule
Channel & Information observed & Primary control \\
\midrule
RFQ & Contract and counterparty; customer side remains hidden & Bid, offer and side-specific displayed quantities \\
FOK & Contract, customer side, price and total quantity & Binary accept/reject subject to full-order risk \\
\bottomrule
\end{tabularx}
\end{table}

\textbf{Research standards.}  The model is deterministic conditional on inputs and
seeds, auditable at every decision, explicit about assumptions, and evaluated using
walk-forward tests, scenario stress tests, ablations and runtime limits.  Bankruptcy
prevention is a hard constraint; PnL optimisation is secondary.

\clearpage

% -----------------------------------------------------------------------------
% Page 3
% -----------------------------------------------------------------------------
\section{Problem definition and market microstructure}
\sectionrule

\subsection{Binary contracts}
A binary option $i$ has legs $(w_{i\ell},U_{i\ell})$, strike $K_i$ and expiry
$T_i$.  Its observable at expiry is
\[
 Z_i(T_i)=\sum_{\ell=1}^{L_i}w_{i\ell}U_{i\ell}(T_i),
\]
and its cash payoff is
\[
 X_i=\1_{\{Z_i(T_i)\ge K_i\}}\in\{0,1\}.
\]
Examples include $\1_{\{\mathrm{FED}_{t+1}\ge1.75\}}$,
$\1_{\{\mathrm{THR}_{t+2}\ge1735\}}$, and
$\1_{\{\mathrm{THR}_{t+1}-\mathrm{AJR}_{t+1}\ge0\}}$.  With zero interest and
no separate risk-premium adjustment, the model fair value is the physical-model
probability
\[
 p_i(t)=\E_t[X_i]=\mathbb{P}_t\!\left(Z_i(T_i)\ge K_i\right).
\]
This is a modelling convention of the simulated market, not a general statement that
real-world derivative prices equal physical probabilities.

\subsection{RFQ protocol}
On an RFQ the exchange requests a two-sided quote
\[
 Q_i=(b_i,n_i^b,a_i,n_i^a),\qquad 0\le b_i<a_i\le1.
\]
The dealer does not know whether the customer will buy or sell.  If the customer
sells, high bids are competitive; if the customer buys, low offers are competitive.
Displayed quantity controls the maximum fill allocated to that side.  Consequently,
an RFQ decision must price both directions and cannot condition directly on the
unrevealed customer side.

\subsection{FOK protocol}
On a fill-or-kill order, the customer provides side $s\in\{\mathrm{BUY},
\mathrm{SELL}\}$, price $P$, and quantity $q$.  The dealer returns only
\texttt{True} or \texttt{False}.  It cannot select a smaller quantity.  If several
dealers accept, the exchange may split the complete customer order according to its
allocation rule; nevertheless each dealer must be capable of carrying all $q$ because
the realised allocation is not guaranteed in advance.

\subsection{Sign and cash conventions}
Let dealer quantity $q_d>0$ denote a purchase and $q_d<0$ a sale.  Worst-case cash
reserved at execution is
\[
 L(q_d,P)=
 \begin{cases}
 q_dP,&q_d>0,\\
 (-q_d)(1-P),&q_d<0.
 \end{cases}
\]
A long contract can lose its purchase price; a short contract can lose one minus its
sale price.  The internal cash ledger deducts this amount immediately.  At expiry it
releases the payoff for a long or the unused collateral for a short.  This convention
mirrors the checker and makes solvency a pathwise property rather than an expectation.

\subsection{Ordered objective}
The strategy solves a lexicographic problem:
\begin{enumerate}
  \item maintain non-negative cash after every trade and settlement, with a buffer;
  \item conditional on survival, maximise realised PnL and the fraction of competitors
  outperformed.
\end{enumerate}
Expected edge never overrides the hard solvency constraint.  A positive-expectation
trade can still be rejected if its worst-case loss or concentration is too large.

\clearpage

% -----------------------------------------------------------------------------
% Page 4
% -----------------------------------------------------------------------------
\section{Stochastic factor model and option valuation}
\sectionrule

\subsection{FED dynamics}
Let $R_t$ denote the policy-rate state on a grid with step $\delta_R=0.25$.  Conditional
on $R_t=r$, the next value is $r+\delta_R$, $r-\delta_R$ (floored at zero), or $r$.
Mean reversion tilts the base transition probabilities:
\[
 \pi_t^{\uparrow}=\clip\!\left(\pi^{\uparrow}+\kappa(R^\star-R_t),0,1\right),
\]
\[
 \pi_t^{\downarrow}=\clip\!\left(\pi^{\downarrow}-\kappa(R^\star-R_t),
 0,1-\pi_t^{\uparrow}\right),
\]
with $R^\star=2$.  A finite-state recursion prices single-leg FED options exactly and
avoids Monte Carlo noise.

\subsection{AJR and THR dynamics}
Let $V_{j,t}$ be company $j\in\{A,H\}$, where $A$ denotes AJR and $H$ denotes THR.
The one-step log return is
\[
 r_{j,t+1}^{\mathrm{company}}
 :=\log\frac{V_{j,t+1}}{V_{j,t}}
 =\mu_j+\beta_{j,R}\Delta R_{t+1}
 +\beta_{j,S}S_{t+1}+\varepsilon_{j,t+1}.
\]
The common sector shock and idiosyncratic residuals satisfy
\[
 S_{t+1}\sim N(0,\sigma_S^2),\qquad
 \varepsilon_{j,t+1}\sim N(0,\sigma_{j,\varepsilon}^2),
\]
independently across the two residuals and independently of $S$.  Therefore
\[
 \Cov(r_{A,t+1},r_{H,t+1}\mid\Delta R_{t+1})
 =\beta_{A,S}\beta_{H,S}\sigma_S^2.
\]
The use of log returns preserves positivity before the simulation's cent rounding and
turns multiplicative growth into an additive Gaussian factor model.  Conditional on a
FED path, a single company value is approximately lognormal.  The log ratio
$\log(V_A/V_H)$ is normal because it is a linear combination of jointly normal log
returns; this permits fast analytic pricing of zero-strike AJR--THR comparisons.

\subsection{Hybrid pricing engine}
The live engine selects the cheapest sufficiently accurate method:
\begin{itemize}
  \item exact dynamic programming for single-leg FED options;
  \item conditional lognormal formulae for single AJR or THR legs;
  \item a conditional normal formula for opposite-sign AJR--THR legs with zero strike;
  \item deterministic Monte Carlo for unsupported combinations.
\end{itemize}
The research configuration uses
\[
 N_{\mathrm{live}}=1000,\qquad N_{\mathrm{theo}}=5000.
\]
The first controls live fallback pricing; the second is reserved for the independent
theoretical-pricing interface.  Each simulation seed is a deterministic function of
parameters, option and state.  Repeated calls therefore return identical probabilities,
which is essential for attribution and debugging.

\subsection{Pricing limitations}
The fitted dynamics are deliberately parsimonious.  They omit jumps, stochastic
volatility, regime changes, endogenous market impact and a separate risk-neutral
measure.  Daily cent rounding also makes analytic lognormal formulae approximate for
some contracts.  These limitations enter execution through model-error and volatility
buffers rather than being hidden.  Production deployment would require independent
calibration validation and monitoring for structural breaks.

\clearpage

% -----------------------------------------------------------------------------
% Page 5
% -----------------------------------------------------------------------------
\section{Rolling calibration and uncertainty control}
\sectionrule

\subsection{History window and exponential weighting}
The calibration consumes only history actually supplied before trading.  It never
assumes that 60 or 200 observations exist.  The maximum window is
$W=60$ observations and the effective half-life is
\[
 H=\min\{20,\max(5,(N-1)/2)\}.
\]
With decay $\lambda=\exp(-\log2/H)$, observation $k$ receives weight proportional to
$\lambda^{N-1-k}$.  The method therefore adapts rapidly under short burn-in while
retaining stability when more data are available.

\subsection{Rate and company estimation}
Weighted up/down indicators estimate the FED transition probabilities.  Regressing the
difference between up and down indicators on $R^\star-R_t$ estimates the mean-reversion
coefficient $\kappa$.  Company log returns are fitted by weighted least squares:
\[
 r_{j,t+1}=\mu_j+\beta_{j,R}\Delta R_{t+1}+e_{j,t+1}.
\]
The weighted residual variances and covariance form
\[
 \Sigma_e=
 \begin{pmatrix}
 \sigma_A^2&\sigma_{AH}\\
 \sigma_{AH}&\sigma_H^2
 \end{pmatrix}.
\]
A one-factor decomposition chooses sector loadings reproducing $\sigma_{AH}$ and
assigns the remaining non-negative variance to idiosyncratic shocks.  Since only the
products $\beta_{j,S}\sigma_S$ are identifiable, the implementation normalises
$\sigma_S=1$.

\subsection{Shrinkage toward a prior}
Short samples are noisy.  With prior sample size $N_0=30$, the empirical weight is
\[
 w_N=\frac{N}{N+N_0}.
\]
Transition probabilities, mean reversion and residual variances are blended with
conservative defaults.  Drifts and rate betas are multiplied by $w_N$.  This reduces
overfitting without imposing a fictitious 200-day history.

\subsection{Estimation-uncertainty spread}
For fair value $p$, the heuristic standard error is
\[
 \operatorname{SE}(p)=
 \sqrt{\frac{\max\{p(1-p),0.01\}}{\max\{N,25\}}}.
\]
The quote adds $a_{\mathrm{model}}\operatorname{SE}(p)$ to its half-spread, with
$a_{\mathrm{model}}=0.25$ in the competitive RFQ configuration.  At $p=0.5$ and
$N=25$, this term is $0.025$; at $N=100$, it falls to $0.0125$.  Thus uncertainty
declines smoothly with evidence, but never disappears at extreme probabilities.

\subsection{Live volatility and markouts}
For each payoff-exposure key, the system records one-step changes in live fair value,
excluding terminal zero/one jumps.  After at least three observations,
\[
 \widehat\sigma_{\mathrm{live}}
 =\frac{n}{n+8}\sqrt{\frac1n\sum_{k=1}^{n}(\Delta p_k)^2},
\]
clipped at $0.06$.  The spread contribution is
$\min\{0.02,0.35\widehat\sigma_{\mathrm{live}}\}$.  Execution markouts are also
recorded by counterparty and payoff key.  A negative direction-adjusted markout means
the market moved in the customer's favour and supplies evidence of adverse selection.

\subsection{Calibration governance}
All fitted parameters, window lengths and priors should be logged at each step.  Alerts
should cover probabilities at constraints, covariance matrices near singularity,
large parameter jumps and systematic forecast miscalibration.  Research changes must
be tested out of sample; hidden-test scores are not themselves a statistically valid
calibration objective.

\clearpage

% -----------------------------------------------------------------------------
% Page 6
% -----------------------------------------------------------------------------
\section{RFQ pricing: reservation price and adaptive spread}
\sectionrule

\subsection{Inventory-aware centre}
Let $X$ be the candidate payoff and let the active portfolio be
$\Pi=\sum_jq_jX_j$.  The relevant covariance is
\[
 C=\Cov(X,\Pi)=\sum_jq_j\Cov(X,X_j).
\]
Dynamic risk aversion is
\[
 \gamma_t=\clip\!\left(
 \gamma_0[1+2D_t+0.75U_t+2L_t],0.01,0.10\right),
\]
where $D_t$ is realised drawdown, $U_t$ is collateral utilisation and $L_t$ is a
recent-negative-PnL fraction.  With $\gamma_0=0.03$, define
\[
 s_t^{\mathrm{inv}}=\clip(\gamma_tC,-0.12,0.12).
\]
Positive covariance with a long portfolio lowers the centre, discouraging more buying
and encouraging sales; negative covariance recognises hedging benefit.

\subsection{Learned flow adjustment}
For each payoff key the system fits
\[
 \Delta p_{t+1}=\alpha+\theta F_t+\eta_{t+1},
\]
shrinks $\theta$ toward zero and clips it to $[-0.03,0.03]$.  Customer buying creates
positive $F_t$; selling creates negative $F_t$; the signal decays by one half at each
step.  The price adjustment is $\clip(\theta F_t,-0.02,0.02)$.  The reservation price
is consequently
\[
 r=\clip\!\left(p-s_t^{\mathrm{inv}}+\clip(\theta F_t,-0.02,0.02),0,1\right).
\]
Inventory and flow move both bid and offer in the same direction; they do not merely
widen the quote.

\subsection{Half-spread decomposition}
The symmetric risk width is
\[
\begin{aligned}
 h=\clip\bigl(&h_{0,t}
 +a_m\operatorname{SE}(p)
 +a_vp(1-p)
 +a_T\min(\sqrt{T},3)\\
 &+a_{\tau}\tau
 +\min(0.02,0.35\widehat\sigma_{\mathrm{live}})
 +h_{\mathrm{repeat}},\ 0.01,h_{\max}\bigr).
\end{aligned}
\]
The recommended research settings are
\[
 h_0=0.01,\quad a_m=0.25,\quad a_v=0.01,\quad
 a_T=0.002,\quad a_{\tau}=1,
\]
with adaptive-base range $[0.01,0.04]$ and final cap $h_{\max}=0.06$.  The repeat
penalty is $\min\{0.015,0.003(m-1)_+\}$ for the $m$th request from the same
counterparty during a step.

Counterparty toxicity blends global and payoff-specific adverse markouts:
\[
 \tau=\clip(0.25\tau_{\mathrm{global}}+0.75\tau_{\mathrm{specific}},0,0.10),
\]
where both components are shrunk toward zero under limited observations.  Toxicity
widens both sides because an RFQ does not reveal which side the customer will choose.

\subsection{Penny-grid construction}
Raw prices are $r-h$ and $r+h$.  The dealer rounds outward:
\[
 b=\max\{0,\min(0.99,\lfloor100(r-h)\rfloor/100)\},
\]
\[
 a=\min\{1,\max(0.01,\lceil100(r+h)\rceil/100)\}.
\]
Outward rounding never silently reduces the intended protection.  For $r=0.5823$ and
$h=0.06$, one obtains $b=0.52$ and $a=0.65$.  A selective competitive rule may cap
$h$ at $0.03$ when $0.2\le p\le0.8$, measured toxicity is at most $0.02$, and
$|q_i|\le5$; this would give approximately $0.55/0.62$.  The override must not be
applied to toxic counterparties or concentrated inventories.

\clearpage

% -----------------------------------------------------------------------------
% Page 7
% -----------------------------------------------------------------------------
\section{RFQ size, tail protection and adaptive behaviour}
\sectionrule

\subsection{Worst loss per displayed contract}
At bid $b$, buying one contract has worst loss $\ell_b=b$.  At offer $a$, selling one
has worst loss $\ell_a=1-a$.  Let
\[
 B_t=f_t^{\mathrm{buffer}}C_0,\qquad F_t=(C_t-B_t)_+,
\]
where $C_0$ is initial capital, $C_t$ is internal free cash and
$f_t^{\mathrm{buffer}}\ge0.15$ rises with drawdown, utilisation and recent loss.  Cash
capacity is $\lfloor F_t/\ell\rfloor$ for positive $\ell$.

The target per-fill risk budget is
\[
 R_t=\min\{0.18C_0,0.30F_t\},
\]
in the aggressive RFQ configuration.  Target quantity is approximately
$\max\{1,\lfloor R_t/\ell\rfloor\}$, scaled downward by toxicity and a dynamic risk
scale.  The final displayed size is the minimum of target size, cash capacity, the
maximum quote size (30), and per-option, underlying and gross-position capacity.

\subsection{Risk-reducing priority}
A sale against an existing long position, or a purchase against an existing short
position, reduces net economic exposure.  The size logic therefore permits enough
quantity to flatten the position when possible.  Nevertheless, the trade still passes
cash capacity because the simulated exchange reserves maximum loss trade by trade and
does not net collateral automatically.

\subsection{Directional tail protection}
Model error is particularly consequential when a customer asks the dealer to take the
unlikely side of an extreme-probability event.  The recommended overlay is
\begin{lstlisting}
if fair_value < 0.20 and current_inventory >= 0:
    bid_quantity = max(1, int(0.40 * bid_quantity))

if fair_value > 0.80 and current_inventory <= 0:
    offer_quantity = max(1, int(0.40 * offer_quantity))
\end{lstlisting}
At low $p$, the dangerous side is buying; at high $p$, the dangerous side is selling.
The opposite side is not reduced because it either takes the statistically favourable
payoff direction or decreases an existing position.

\subsection{Adaptive base spread}
After each step the base half-spread changes by
\[
 \Delta h_{0,t}=0.04(-\overline m_t)
 +0.003(\overline f_t-0.30)+0.005L_t,
\]
where $\overline m_t$ is average direction-adjusted information markout and
$\overline f_t$ is average fill rate.  Adverse markouts, excessive fills and negative
PnL widen the base.  Favourable markouts and very low participation narrow it.  The
update is deliberately slow and clipped to avoid instability.

\subsection{RFQ algorithm}
\begin{enumerate}
  \item Price the contract and retrieve portfolio covariance, flow and toxicity.
  \item Compute reservation price $r$ and half-spread $h$.
  \item Apply the central-probability competitive cap only under low toxicity and
  small inventory.
  \item Round bid downward and offer upward to whole pennies.
  \item Compute side-specific quantities from worst loss, free cash and hard limits.
  \item Reduce dangerous tail-side quantity; preserve risk-reducing capacity.
  \item If one side has zero capacity, quote the zero-worst-loss boundary price with
  minimum legal size: bid zero or offer one.
\end{enumerate}

The design is intentionally asymmetric: aggressiveness is concentrated where the
estimated event probability is central and risk can be diversified; extreme tail
exposure is acquired slowly.

\clearpage

% -----------------------------------------------------------------------------
% Page 8
% -----------------------------------------------------------------------------
\section{FOK acceptance and adverse-selection protection}
\sectionrule

\subsection{Raw edge}
For a customer BUY, the dealer sells at $P$ and has raw edge
\[
 e_{\mathrm{raw}}=P-r,\qquad L_{\max}=q(1-P).
\]
For a customer SELL, the dealer buys and has
\[
 e_{\mathrm{raw}}=r-P,\qquad L_{\max}=qP.
\]
Using the reservation price rather than unadjusted fair value incorporates the marginal
effect on the existing portfolio.

\subsection{Information hurdle}
The base required edge is
\[
\begin{aligned}
 e_{\mathrm{req}}=\clip\bigl(&e_0
 +a_q\min(\sqrt q,4)
 +0.75\tau
 +a_m^{F}\operatorname{SE}(p)\\
 &+\min(0.015,0.25\widehat\sigma_{\mathrm{live}})
 +h_{\mathrm{repeat}},\ e_0,e_{\max}\bigr).
\end{aligned}
\]
The conservative settings are
\[
 e_0=0.015,\quad a_q=0.010,\quad a_m^F=0.40,\quad e_{\max}=0.20.
\]
Increasing $e_{\max}$ alone does not widen the hurdle; it only allows the sum of
other penalties to exceed the previous cap.

Large orders receive an additional concentration charge
\[
 e_{\mathrm{conc}}=\min\{0.15,\gamma_tq\,p(1-p)\},
\]
and short maturities receive
\[
 e_{\mathrm{short}}=\min\left\{0.03,
 \frac{0.03}{\max(T,1)}\right\}.
\]
The economic acceptance test is
\[
 e_{\mathrm{raw}}\ge e_{\mathrm{req}}+e_{\mathrm{conc}}+e_{\mathrm{short}}.
\]
This treats large one-day FOKs as potentially informed even when the counterparty has
no prior toxicity history.

\subsection{Solvency and concentration gates}
Let $B_t$ be the dynamic cash buffer.  Affordability requires
\[
 C_t-L_{\max}\ge B_t.
\]
For a risk-increasing FOK, the separate loss budget is
\[
 L_{\mathrm{budget}}
 =\min\{(C_t-B_t)_+,0.08C_0\}.
\]
A fully risk-reducing order may use up to $0.12C_0$, but still must pass affordability
and hard position limits.  The final rule is
\[
 \boxed{\mathrm{accept}=\mathrm{attractive}\wedge\mathrm{affordable}
 \wedge\mathrm{within\ limits}\wedge\mathrm{within\ loss\ budget}.}
\]

\subsection{Interpretation of observed behaviour}
Consider selling 17 contracts at $0.78$.  The maximum loss is
$17(1-0.78)=3.74$.  With $C_0=40$ and an eight-percent risk-increasing budget, the
limit is $3.20$, so the order is rejected even before considering information risk.
Selling two contracts at $0.99$ risks only $0.02$ and may remain acceptable when edge
is strong.  This illustrates the desired nonlinearity: contract count alone is not
risk, and notional price alone is not risk; the relevant quantity is payoff-contingent
maximum loss combined with information quality.

\subsection{FOK implementation skeleton}
\begin{lstlisting}
accepted = (
    raw_edge >= required_edge
                + concentration_charge
                + short_tenor_penalty
    and cash_balance - maximum_loss >= buffer
    and position_capacity >= order.quantity
    and maximum_loss <= loss_budget
)
return accepted
\end{lstlisting}

\clearpage

% -----------------------------------------------------------------------------
% Page 9
% -----------------------------------------------------------------------------
\section{Solvency, portfolio controls and runtime architecture}
\sectionrule

\subsection{Exact maximum-loss ledger}
Every fill changes internal cash immediately.  For dealer purchase quantity $q>0$,
\[
 C_{t+}=C_t-qP.
\]
For dealer sale quantity $-q<0$,
\[
 C_{t+}=C_t-q(1-P).
\]
At expiry with payoff $X\in\{0,1\}$, a long receives $qX$ and a short recovers
$q(1-X)$.  Settlement can only increase reserved cash, whereas execution reduces it.
The hard affordability gate is therefore applied before every voluntary FOK and is
embedded in RFQ displayed-size capacity.

\subsection{Layered risk controls}
\begin{table}[h]
\centering\small
\begin{tabularx}{\textwidth}{>{\raggedright\arraybackslash\bfseries}p{3.0cm}p{2.6cm}X}
\toprule
Control & Research setting & Function \\
\midrule
Cash buffer & $\ge15\%$ of initial capital & Preserves liquidity under consecutive one-sided fills \\
Per-option cap & 200 contracts & Emergency protection against single-event concentration \\
Underlying cap & 300 contracts & Limits common AJR, THR or FED factor exposure \\
Gross cap & 500 contracts & Bounds total operational exposure \\
Maximum quote size & 30 contracts & Limits any individual RFQ allocation \\
FOK loss budget & $8\%$ risk-increasing & Rejects large all-or-nothing downside \\
Dynamic risk scale & $[0.40,1]$ & Reduces displayed size after loss or high utilisation \\
\bottomrule
\end{tabularx}
\end{table}

Contract caps are backstops, not substitutes for price-sensitive cash control.  Buying
25 contracts at $0.01$ risks $0.25$, while buying 25 at $0.99$ risks $24.75$; a single
hard count treats economically different positions as identical.

\subsection{Fast state-level covariance}
The original implementation repeatedly simulated portfolio covariance after each
fill, although the underlying state and fitted parameters had not changed.  The fast
architecture generates $N_{\mathrm{cov}}=500$ common market paths once per state and
stores each option's sampled binary payoff as an integer bitset.  If $B_X$ and $B_Y$
are payoff bitsets, then
\[
 \widehat{\E}[XY]=\frac{\operatorname{bitcount}(B_X\mathbin{\&}B_Y)}{N_{\mathrm{cov}}}.
\]
Pair covariance is obtained from this joint mean and the two bit counts, and portfolio
covariance is a fast weighted sum.  A trade changes $q_j$ but not the simulated joint
distribution, so it must not invalidate scenario or fair-value caches.  Caches are
cleared only when the market state or fitted model changes.

Reducing covariance paths from 1500 to 500 triples the speed of scenario generation
but increases Monte Carlo standard error by $\sqrt3$.  The cache changes the more
important dimension: simulation cost is paid once per market state rather than once
per order or fill.  Live prices are similarly cached by immutable option within the
state.  The separate theoretical test retains its own deterministic sample count.

\subsection{Operational safeguards}
Production-style controls should include bounded logging, latency counters, cache-hit
rates, exception isolation, invariant checks ($0\le b<a\le1$ and positive sizes), and
a deterministic replay facility.  Debug printing must be disabled in scored or live
sessions because unbounded output can itself cause operational failure.

\clearpage

% -----------------------------------------------------------------------------
% Page 10
% -----------------------------------------------------------------------------
\section{Evaluation, governance and conclusions}
\sectionrule

\subsection{Validation protocol}
A credible evaluation separates parameter fitting from performance assessment.
Recommended evidence includes:
\begin{enumerate}
  \item \textbf{Walk-forward tests:} fit only on data available before each session;
  never use future outcomes to tune a contemporaneous quote.
  \item \textbf{Competitor diversity:} evaluate against fixed-width, inactive,
  inventory-aware and situational strategies, with randomised order sequences.
  \item \textbf{Ablations:} remove covariance, toxicity, flow, adaptive spread and FOK
  penalties one at a time to identify genuine contribution.
  \item \textbf{Stress tests:} long one-sided RFQ flow, repeated toxic counterparties,
  correlated expiries, probabilities near zero/one, small starting capital, and maximum
  order sizes.
  \item \textbf{Runtime tests:} record decision latency, simulation calls and total
  session time separately from PnL.
\end{enumerate}

Primary metrics are bankruptcy frequency, minimum cash, maximum drawdown, realised
PnL, rank, fill rate, spread capture, one-step markout, tail loss, exposure by factor,
FOK acceptance by size/tenor, and runtime percentiles.  Passing a runtime limit only
means the strategy completed; it does not create a PnL score.  Conversely, a
non-bankrupt strategy can receive a low score if it beats no competitor.

\subsection{Parameter governance}
The displayed constants are research starting points, not universal optima.  Each
change should have an economic hypothesis and an acceptance criterion.  For example,
tightening the central RFQ spread should increase competitive fills without worsening
adverse markouts beyond a pre-specified threshold.  Reducing the FOK loss fraction
should lower tail loss without eliminating all high-edge small orders.  Multiple
parameters should not be changed simultaneously when attribution matters.

Model versions should record code hash, parameter set, random-seed convention, training
window, test set and exception logs.  An independent reviewer should challenge model
assumptions, probability calibration, covariance estimation, collateral accounting and
the possibility that counterparties condition orders on information absent from the
model.

\subsection{Conclusion}
The proposed market maker is organised around three principles.  First, price and risk
are distinct: fair value estimates the event probability, the reservation price manages
directional exposure, and the spread compensates for uncertainty and information risk.
Second, execution protocols require different behaviour: ordinary central-probability
RFQs may justify aggressive two-sided competition, while revealed-side FOKs require
size, concentration and short-tenor penalties.  Third, solvency is pathwise: every
accepted order must be affordable under its worst payoff, irrespective of expected
edge.  State-level caching makes this discipline computationally feasible without
removing the economic controls.

The framework remains a research prototype.  Its next development priorities are
probability-calibration diagnostics, explicit regime detection, uncertainty intervals
for portfolio covariance, and walk-forward optimisation of RFQ aggressiveness subject
to a hard bankruptcy constraint.

\vspace{2pt}
{\small
\begin{thebibliography}{9}
\setlength{\itemsep}{0pt}
\bibitem{avellaneda}
M.~Avellaneda and S.~Stoikov, ``High-frequency trading in a limit order book,''
\emph{Quantitative Finance}, 8(3), 217--224, 2008.
\bibitem{glosten}
L.~R.~Glosten and P.~R.~Milgrom, ``Bid, ask and transaction prices in a specialist
market with heterogeneously informed traders,'' \emph{Journal of Financial Economics},
14(1), 71--100, 1985.
\bibitem{ho-stoll}
T.~Ho and H.~R.~Stoll, ``Optimal dealer pricing under transactions and return
uncertainty,'' \emph{Journal of Financial Economics}, 9(1), 47--73, 1981.
\bibitem{glasserman}
P.~Glasserman, \emph{Monte Carlo Methods in Financial Engineering}, Springer, 2004.
\bibitem{gueant}
O.~Gu\'eant, C.-A.~Lehalle and J.~Fernandez-Tapia, ``Dealing with the inventory risk,''
\emph{Mathematics and Financial Economics}, 7, 477--507, 2013.
\end{thebibliography}
}

\end{document}
